\chapter{混音进阶实验}
\label{chap:advanced-lab}

\classquote{井蛙不可以语于海者，拘于虚也；夏虫不可以语于冰者，笃于时也。}{《庄子·秋水》}{简注：井蛙被空间所限，夏虫被时令所限；进阶实验就是主动走出熟悉的参数与听感范围。}

前 10 章把夜半的信号路径、参数契约与基本判断讲完了。这一章不重讲基础，而是把常见进阶技巧拆成可以在夜半里动手验证的实验：每一项先说“要解决什么听感问题”，再说“当前版本能做什么、不能做什么”，最后给一个可复现的练耳/操作步骤。不能在本机架上完成的技巧会明确标出边界，不拿通用 DAW 的按钮冒充夜半的入口。

\section{电平预算与平台响度}

\term{电平预算}（level budget）是在整条链上预先分配峰值的做法：音源、设备链、总线、导出各留多少余量。它和“把推子推到多响”不是一回事。一个稳妥的起点是：单轨峰值不超过 $-6$ dBFS，主总线在未限制导出时留 6 dB 余量；需要更响时，优先检查编曲密度和动态，而不是把每一级都顶满。

\term{平台响度归一化}（platform loudness normalization）是流媒体平台在播放时把不同作品拉到相近响度的机制。它不要求所有歌都做成同一个 LUFS：平台会按自己的目标调整，过度压扁的作品反而会在归一化后失去动态。工程里可以把“交付目标”设成一个参考值，但不要把“平台会归一化”理解成“母带可以随便超”。

\begin{table}[htbp]
\centering
\small
\begin{tabular}{@{}p{0.24\linewidth}p{0.22\linewidth}p{0.44\linewidth}@{}}
\hline
口径 & 常见参考 & 说明 \\
\hline
流媒体立体声 & $-14$ LUFS 左右 & 常见平台目标区间；不是所有平台都相同 \\
播客/对白 & $-16$ LUFS 左右 & 以语言可懂度为先，音乐作品不套用 \\
EBU R128 & $-23$ LUFS & 广播交付口径；不是流媒体默认 \\
真峰值上限 & $-1$ dBTP & 给编码/重采样留余量；不是“必须正好 $-1$” \\
未限制余量 & 约 6 dB & 给限制器与后续处理留空间 \\
\hline
\end{tabular}
\caption{常见响度参考；目标数字是工程选择，不是规范强制}
\label{tab:advanced-loudness-targets}
\end{table}

夜半的边界必须说清楚：混音台主表上的“LUFS”是平滑 RMS 减 0.691 dB，不能当作 BS.1770 门限积分；控制室的“积分 LUFS”和“dBTP”两格在当前构建里没有运行时更新；唯一可交付的读数是导出后的实测行。要在导出时追 $-14$ LUFS 或 $-1$ dBTP，只能走 CLI/MCP 的响度归一化/真峰值上限；GUI 导出对话框没有这些输入框。不要按“在控制室看到 $-14$ 就达标”写操作步骤。

\subsection{实验 1：先测，再决定}

\begin{enumerate}
  \item 选一段 8–12 秒的副歌，保持所有轨推子不动，导出一版关闭母线限制器的 WAV。
  \item 读导出实测行的真峰值、积分 LUFS 与 LRA；把这三个数抄下来，不要凭耳朵猜“大概 -14”。
  \item 再打开母线限制器导出一次，比较两次的真峰值与 LRA。若 LRA 明显变小而积分 LUFS 只涨了一点，说明限制器在压动态，不是在“提高响度”。
  \item 回到工程，把最响的三条轨各降 1–2 dB，再重复一次；观察真峰值是否下降、整体听感是否更松。记录“降源”和“拉总线”两种做法的差别。
\end{enumerate}

\section{扫频、动态 EQ 与多段思路}

\term{扫频}（frequency sweep）是临时把一段 EQ 提升到很夸张的位置，再缓慢移动中心频率，用“哪里最难受”来定位问题。它只是探针；找到位置后要把提升拉回，改成小幅衰减。夜半的三段 EQ 没有独立 Q 旋钮，曲线拖动只改增益与中心频率，所以探针会比专用 EQ 宽；判断时以整体回听为准。

\term{动态 EQ}（dynamic EQ）是“只在某段超过阈值时才改变该段增益”的处理器。它和静态 EQ 的差别不是曲线形状，而是触发条件：静态 EQ 会一直削，动态 EQ 只在需要时削。夜半的通道条里确实有动态 EQ 级（\texttt{deq\_freq}、\texttt{deq\_threshold\_db}、\texttt{deq\_cut\_db}、\texttt{deq\_q}、\texttt{deq\_ratio}），位置在 EQ 之后、滤波器之前；但设备卡当前没有这五项的独立旋钮，它们只能经 MCP 的 \texttt{DeviceParam} 泳道写入，不能靠拖动三段 EQ 的曲线来调用。三段塑形 EQ 是静态 biquad，动态 EQ 是另一级，两者不要混写。

\term{多段压缩}（multiband compression）是把频谱分成几段，各段独立压缩，再合成。它能只压低频的峰值而不碰高频，但会引入分频相位与段间交互。夜半当前没有多段压缩器件；设备机架只有单段压缩器。要做“只压低频”的实验，只能用静态 EQ 先把低频修到可接受，再用单段压缩整体控制，或把低频单独送到一条辅助轨处理（如果工程允许拆轨）；不要把单段压缩器说成多段。

\subsection{实验 2：扫频找“浑浊”与“刺耳”}

\begin{enumerate}
  \item 选一条人声或贝斯，独奏，装三段 EQ；把中频增益临时推到 $+9$ dB。
  \item 从 200 Hz 缓慢向 2 kHz 移动，停在“最闷/最纸箱”的位置；把增益拉回 0，再向下衰减 2–3 dB。
  \item 再把中频探针移到 2–6 kHz，找“最刺/最齿音”的位置；同样拉回 0，做 1–2 dB 衰减。不要同时大幅提升高频去“补亮度”。
  \item 回到全混，比较“独奏好听”和“整体好听”的差别；如果整体变薄，把衰减减半。
\end{enumerate}

\section{侧链、门限与并行压缩}

\term{侧链压缩}（sidechain compression）是用一条轨的电平去触发另一条轨的压缩。典型用法是 kick 触发 bass，让低频在 kick 落点时短暂退让；或人声触发音乐总线，让人声出来时伴奏自动让路。夜半的设备机架里有“动态、侧链与瞬态控制”面板，第一行“侧链键源”可以点击循环：它把展示轨的键源依次切到无、其他轨；选中一条源轨后，路由里会有一条 \texttt{Sidechain} 边，通道条压缩器的检波器可以在有键时读键信号。当前面板没有独立的侧链阈值/深度旋钮，键源选择之外的参数仍按压缩器契约写入；不要把它说成“没有侧链”，也不要把普通阈值旋钮说成侧链深度。

\term{门限}（gate）是在阈值以下关闭或大幅衰减信号，用来去底噪、串音和不需要的尾音；\term{扩展器}（expander）则在阈值以下进一步压低，目的是加大电平差，不是静音。两者都作用于阈值以下，和压缩器作用于阈值以上相反。夜半的 DSP 里有独立的 \texttt{gate.rs}（Gate / Expander / Ducker 三种模式），但当前没有接进设备链、也没有机架或面板入口；不能把压缩器比率调高来代替门限：压缩器不会在静音段把信号关掉。

\term{并行压缩}（parallel compression）是把同一信号的干声与重压版本并联，再按比例混合：保留瞬态与细节，同时把持续段压得更实。夜半可以做一条近似路径：建一条辅助返回轨，在返回轨上装压缩器，源轨通过发送把信号送进去，返回轨再回主总线；干声仍走原来的轨。发送/返回的入口见第 5 章，压缩器参数见第 4 章。注意这不是“插件里的并联开关”，而是两条信号路径的汇合；返回轨的压缩器会同时处理所有送进来的轨。

\subsection{实验 3：用发送/返回做一次并行压缩}

\begin{enumerate}
  \item 选一条鼓组或人声轨，建一条辅助返回轨；在返回轨上装压缩器。
  \item 把返回轨压缩器的阈值压低、比率调到 4:1–8:1，让返回信号明显更“实”；先不要管返回轨推子。
  \item 源轨发送增益从 $-20$ dB 开始，慢慢推到能听见“实”的质感但不盖过干声；通常只需要很少的量。
  \item 旁通返回轨压缩器，比较干声与并联结果；如果差别只是“更响”，先把返回轨推子降回去，再判断质感是否真的变好。
\end{enumerate}

\section{延迟、Haas 与宽度}

\term{预延迟}（pre-delay）是湿声相对干声整体后推的时间。它让直达声先站住，再进空间；太短会把干声和尾巴糊在一起，太长会让声音“离墙太远”。夜半的 Freeverb 有 \texttt{reverb\_predelay} 参数，但设备卡当前没有可拧的旋钮；它只能经工程/MCP 写入，不能按“界面上拖一下”来写。

\term{Haas 延迟}（Haas delay）是把一个声音复制到另一侧、延迟 1–35 ms，让听者把先到的那个定位成声源。它常被误当成“加宽”，但单声道下会变成梳状滤波；真正的宽度检查必须回到 MONO。夜半当前没有独立的 Haas 延迟器件，也没有“复制到另一侧并延迟”的界面动作；不能把轨道声相或混响说成 Haas。

宽度在夜半不是轨道推子旁的旋钮，而是通道条的 \texttt{strip\_width}（M/S 缩放）；\texttt{mono\_below\_hz} 可以在低频做 mono-maker。两者都只在挂了通道条时存在，当前设备卡没有宽度/mono-maker 旋钮：要改数值，走 MCP 的 \texttt{DeviceParam} 泳道，按名写 \texttt{strip\_width} / \texttt{mono\_below\_hz}，或直接改工程里的设备参数。单声道兼容要用控制室的 MONO/SIDE 监听检查：MONO 下消失的必留声部，说明左右相消；SIDE 下该留的中心干声消失，说明中频被差信号吃掉。监听开关只改监听，不进导出。

\subsection{实验 4：宽度不是“越宽越好”}

\begin{enumerate}
  \item 选一条有立体声宽度的轨（铺底、和声或鼓组），先用 MONO/SIDE 听它当前的宽度是否在单声道下塌陷；把结果记下来。
  \item 若能用 MCP，按通道条契约把 \texttt{strip\_width} 从 1.0 写到 1.4，再写到 0.6；每次回到全混听一遍，并重复 MONO 检查。不能用 MCP 时不要假装设备卡上有宽度旋钮，改做第 3 步。
  \item 打开控制室 MONO，比较不同已有素材的单声道结果；记录哪一种在 MONO 下损失最少。
  \item 打开 SIDE，听差信号里还剩什么；如果中心人声/贝斯在 SIDE 里仍然很响，说明它的中心成分过多，回到 MONO 再确认。
  \item 不要把 MONO 打开时的监听电平当作导出电平；MONO/SIDE 只是检查开关。
\end{enumerate}

\section{染色、并行饱和与监听环境}

\term{谐波失真}（harmonic distortion）是非线性处理产生的整数倍频率；\term{饱和}（saturation）是让非线性在接近满刻度时自然压缩，产生谐波但不完全削平。夜半的磁带饱和是链级器件：\texttt{tape\_drive} 与 \texttt{tape\_mix} 同时大于 0 才建；它是串联在链里的，不是一个可以只送湿声的并行效果。

夜半当前没有独立的并行饱和路由：不能用发送/返回把“只经过磁带级的湿声”送回主总线，也不能在设备卡上把磁带级拆成干/湿两条独立通路。想做并行饱和，只能把同一素材放到另一条轨（若工程支持复制轨），在副本上做不同的染色，再平衡；这会增加一条真实轨，必须确认它不会破坏原来的编曲与导出。不要把“磁带级有 \texttt{tape\_mix}”误读成“并联饱和”：\texttt{tape\_mix} 是链内的干湿比，返回的仍是同一条链。

\term{参考监听}（reference monitoring）不是“找一首好听的歌当背景”，而是用固定的参考曲、固定的监听音量和固定的检查顺序，判断自己的混音在相同条件下的偏差。夜半的参考曲应放在单独音频轨上，旁通该轨效果，只做 A/B；它不进入导出。监听环境至少要检查三个条件：正常音量、DIM 衰减后的音量、MONO/SIDE 下的单声道与差信号。

\subsection{实验 5：染色 A/B 与参考监听}

\begin{enumerate}
  \item 选一条鼓或贝斯，把磁带驱动的 \texttt{tape\_drive} 从 0 调到 0.25，再把 \texttt{tape\_mix} 从 0 调到 0.5；用 \texttt{tape\_trim\_db} 把整体电平补回 A/B 前的样子。
  \item 在相同监听音量下做 A/B：先听瞬态是否变软，再听低频是否更“厚”，最后回到全混判断是否只是变响。
  \item 导入一段参考曲到单独音频轨，旁通其效果；用同样的监听音量与 MONO/SIDE 顺序比较，不要用参考曲的响度直接判断自己的电平。
  \item 记录三条：瞬态、低频、噪声/齿音。没有变化的项要写“无变化”，不要为了写满而编感受。
\end{enumerate}

\section{故障排查：浑浊、刺耳、相位与低频打架}

进阶技巧最后都要落到“哪里不对、先改哪一层”。下面四类问题在夜半里各有可验证的检查顺序；不要一上来就加 EQ 或加压缩。

\begin{table}[htbp]
\centering
\small
\begin{tabular}{@{}p{0.16\linewidth}p{0.34\linewidth}p{0.40\linewidth}@{}}
\hline
症状 & 先查 & 常见修法 \\
\hline
浑浊 & 250–500 Hz 是否多轨堆积；低频是否过量；混响尾巴是否太长 & 先做窄衰减；让 bass/kick 错开；缩短 \texttt{reverb\_decay\_s} 或降低 wet \\
刺耳 & 2–5 kHz 是否有齿音/镲片叠加；高频架是否整体抬高 & 先在中频做 1–3 dB 衰减；检查是否每轨都在抬高频；需要动态 EQ 时，只能走通道条 \texttt{deq\_*}（MCP 写入） \\
相位抵消 & MONO 下是否某声部消失；SIDE 是否异常突出；是否有延迟/极性反转 & 查 \texttt{sample\_delay\_invert} 与并联路径；MONO/SIDE 复听；不要用宽度硬盖 \\
低频打架 & kick 与 bass 的基频/低频包络是否重叠；mono-maker 是否开 & 先调编曲与声相；用低架小幅衰减；让 kick 与 bass 在不同频段/时间让位 \\
\hline
\end{tabular}
\caption{四类常见问题的检查顺序}
\label{tab:advanced-troubleshooting}
\end{table}

\term{相位抵消}（phase cancellation）是两路相同信号在某个频率上反相或延迟后相加变小。它不一定“听起来像坏了”：有时只是某件乐器变薄。检查相位不要只看波形；用 MONO/SIDE、极性反转与延迟参数逐项复听，确认是并联路径还是单轨内部的问题。

\term{低频掩蔽}（low-frequency masking）是低频能量大时，人耳更难听到同时出现的中高频细节。低频“打架”往往不是 EQ 能单独解决的：先检查 kick 与 bass 的时间关系、音头是否撞在一起、声相是否都挤在中间，再决定衰减或让位。

\subsection{实验 6：用 MONO/SIDE 做故障定位}

\begin{enumerate}
  \item 选一段有鼓、贝斯、人声的 8 小节，先听正常立体声；记下“浑浊/刺耳/某声部消失”的具体位置。
  \item 打开 MONO，逐一独奏鼓、贝斯、人声，确认哪些声部在单声道下损失最大；不要在 MONO 下判断宽度。
  \item 打开 SIDE，确认差信号里是否有不该突出的中心干声；如果有，回到 MONO 检查该轨的声相与宽度。
  \item 对怀疑的轨做一次 1–2 dB 的低架衰减或 1–2 dB 的中频衰减，再回到全混；一次只改一个变量。
\end{enumerate}

\section{练习}


\begin{exercise}
本章的实验都只改一个变量，并在每一步回到全混。做完后把结果写进工程备注，不要只凭记忆。

\begin{enumerate}
  \item 用导出实测行记录一版未限制 WAV 的真峰值、积分 LUFS 与 LRA；再记录一版开限制器的结果，解释两者差异。
  \item 用三段 EQ 的扫频法找一处 250–500 Hz 的堆积与一处 2–5 kHz 的刺耳，各做 1–3 dB 衰减，并说明为什么不是提升。
  \item 用辅助返回轨与压缩器做一次并行压缩，比较干声、并联与“只是更响”的差别；记录发送增益。
  \item 选一条有宽度的轨，比较 \texttt{strip\_width}=0.6/1.0/1.4 在 MONO 下的损失；不要用 MONO 监听电平判断导出。
  \item 做一次磁带饱和 A/B，必须用 \texttt{tape\_trim\_db} 把电平补回；记录瞬态、低频、噪声三项。
  \item 用参考曲、正常音量、DIM、MONO/SIDE 四个条件检查同一段混音；列出至少两条偏差。
  \item 对“浑浊、刺耳、相位抵消、低频打架”各写一条检查顺序；每条只改一个变量。
  \item 把以上七项合成一份 8 小节复混报告：每个决定写“听感—证据—改动—回听结果”，不要只写“更好听了”。
\end{enumerate}
\end{exercise}

\section{本章与帮助主题的映射}

\begin{table}[htbp]
\centering
\small
\begin{tabular}{@{}p{0.42\linewidth}p{0.48\linewidth}@{}}
\hline
小节 & 应用内帮助锚点 \\
\hline
电平预算与平台响度 & \texttt{mixing-level-budget} \\
扫频、动态 EQ 与多段思路 & \texttt{mixing-eq-advanced} \\
侧链、门限与并行压缩 & \texttt{mixing-dynamics-advanced} \\
延迟、Haas 与宽度 & \texttt{mixing-space-advanced} \\
染色、并行饱和与监听环境 & \texttt{mixing-coloration-advanced} \\
故障排查：浑浊、刺耳、相位与低频打架 & \texttt{mixing-troubleshooting} \\
\hline
\end{tabular}
\caption{第十一章帮助主题映射}
\end{table}

\section*{延伸阅读}
\begin{itemize}[leftmargin=1.5em]
\item \url{https://www.soundonsound.com/reviews}——Sound On Sound 评测库：读进阶工具与实验性工作流的评测，先判断“值不值得试”。
\item \url{https://www.freesound.org/}——Freesound 素材库：找可合法试听的素材做实验，补足“换一种声音再听同一问题”的练习。
\end{itemize}
